\documentclass[11pt]{article}
\usepackage[T1]{fontenc}
\usepackage{lmodern}
\usepackage{amsmath,amssymb,amsthm,mathtools}
\usepackage[margin=1in]{geometry}
\usepackage{microtype}
\usepackage{booktabs,tabularx,graphicx}
\usepackage[shortlabels]{enumitem}
\usepackage{needspace}
\usepackage{xurl}
\usepackage[hidelinks]{hyperref}
\hypersetup{pdftitle={Local root conditioning and companion eigenvectors of Gaussian polynomials},
 pdfauthor={Etna Lindy, Martin Lotz, Vanni Noferini}}
\setlength{\emergencystretch}{2em}

\newcommand{\C}{\mathbb C}
\newcommand{\R}{\mathbb R}
\newcommand{\E}{\mathbb E}
\newcommand{\Prob}{\mathbb P}
\newcommand{\Estar}{\mathbb E_*}
\newcommand{\Pstar}{\mathbb P_*}
\newcommand{\CN}{\mathcal{CN}}
\newcommand{\B}{\mathcal B}
\newcommand{\dd}{\,\mathrm d}
\DeclareMathOperator{\diag}{diag}
\DeclareMathOperator{\tr}{tr}
\DeclareMathOperator{\adj}{adj}
\DeclareMathOperator{\Var}{Var}
\DeclareMathOperator{\dist}{dist}
\DeclareMathOperator{\spanop}{span}
\DeclareMathOperator{\rank}{rank}
\DeclareMathOperator{\Cov}{Cov}
\newtheorem{theorem}{Theorem}[section]
\newtheorem{proposition}[theorem]{Proposition}
\newtheorem{lemma}[theorem]{Lemma}
\newtheorem{corollary}[theorem]{Corollary}
\theoremstyle{definition}
\newtheorem{definition}[theorem]{Definition}
\newtheorem{example}[theorem]{Example}
\theoremstyle{remark}
\newtheorem{remark}[theorem]{Remark}

\title{Local root conditioning and companion eigenvectors\\
of Gaussian polynomials}
\author{Etna Lindy, Martin Lotz, Vanni Noferini}
\date{Draft --- 12 September 2026}

\begin{document}
\maketitle
\begin{abstract}
For a polynomial of degree $n$ with independent standard complex Gaussian
coefficients, we prove a degree-independent second-power tail for the
eigenvector condition number of its companion matrix at a uniformly selected
root. Its moments are uniformly bounded precisely below order two, while its
second moment is $2\log n+O(1)$. The expected logarithm of the relative
companion-to-coefficient condition ratio is $\log n+O(1)$. The proof separates
local root sensitivity from interpolation in the companion basis. For Gaussian
systems, marked Kac--Rice and a kernel metric reduce the normalized Jacobian
to an induced Ginibre law, giving a sharp fourth-power conditioning tail.
Exact box root laws and a tridiagonal Gram matrix then yield the companion
results. We specify the perturbation spaces and output metrics throughout;
the comparison concerns local conditioning, rather than stability of a
complete algorithm.
\end{abstract}

\section{Introduction}\label{sec:introduction}

Computing the roots of a polynomial by solving a companion eigenvalue problem
changes the perturbations to which the answer is exposed. A polynomial root
responds to errors in its coefficients; the unstructured eigenvalue condition
permits errors in every entry of the companion matrix. We ask how the conditioning of these two problems
compares for Gaussian input. Our main result gives a sharp probabilistic
description of companion eigenvalue conditioning and a relative comparison
with the original root problem.

Average-condition theorems for matrix polynomials use their coefficient
matrices as input. Applying such a theorem after linearization also requires
its probability law and perturbation norm to match those of the induced
companion problem. Section~\ref{sec:intro-previous} makes this distinction
explicit for the results of Armentano--Beltr\'an.

Consider
\[
 p(x)=\sum_{k=0}^n a_kx^k,
 \qquad a_0,\ldots,a_n\ \text{independent standard complex Gaussians}.
\]
Here each coefficient has independent real and imaginary parts of variance
$1/2$. Let $C_p$ be the $n\times n$ monic companion matrix, with ones on its
subdiagonal and last column $-(a_0,\ldots,a_{n-1})^T/a_n$. First sample $p$,
then choose one of its $n$ roots $Z$ uniformly; write $\Pstar,\Estar$ for this
joint law. Almost surely the roots are simple and nonzero.

For a root $z$, let $r,w$ be right and left eigenvectors of $C_p$ at $z$:
$C_pr=zr$ and $w^*C_p=zw^*$. The quantity
\begin{equation}\label{eq:intro-overlap}
 \Lambda(p,z)=\frac{\|w\|_2\|r\|_2}{|w^*r|}
\end{equation}
is the absolute eigenvalue condition number for arbitrary matrix perturbations
in spectral norm. It is often called the eigenvector overlap. We compare
relative condition numbers using Euclidean norm on polynomial coefficients,
spectral norm on matrices, and relative error in the same scalar root. Their
precise definitions are given in Section~\ref{sec:setup}.

\Needspace{17\baselineskip}
\begin{theorem}[Main result: conditioning of Gaussian companions]\label{thm:intro-main}
There is an absolute constant $C$ such that the following hold.
\begin{enumerate}[(i)]
 \item For every degree $n\ge1$ and every $t>0$,
 \begin{equation}\label{eq:intro-tail}
   \Pstar\{\Lambda(p,Z)>t\}\le\min\{1,Ct^{-2}\}.
 \end{equation}
 Moreover,
 \begin{equation}\label{eq:intro-second}
   \Estar\Lambda(p,Z)^2=2\log n+O(1)
   \qquad(n\to\infty).
 \end{equation}
 Consequently the moments $\Estar\Lambda^q$ are bounded uniformly in degree
 precisely for $0<q<2$, and $\Estar\log\Lambda=O(1)$.
 \item The relative eigenvalue and coefficient condition numbers satisfy
 \begin{equation}\label{eq:intro-relative}
   \Estar\log
   \frac{\kappa_{\mathrm{eig,rel}}(C_p,Z)}
        {\kappa_{\mathrm{coeff,rel}}(p,Z)}
   =\log n+O(1)
   \qquad(n\to\infty).
 \end{equation}
\end{enumerate}
The constants in the $O(1)$ terms are independent of $n$.
\end{theorem}

The first assertion controls the expected fraction of roots with a large
eigenvalue condition number, uniformly in the degree. Its exponent is sharp:
any larger exponent with a degree-independent constant would bound the second
moment, contradicting \eqref{eq:intro-second}. The second assertion says that
the geometric mean of the relative companion-to-coefficient condition ratio
is comparable to $n$. These statements fit together because the coefficient
root condition has its own degree scale. They concern a uniformly selected
root and two specified local problems, rather than a bound on all errors in
a complete rootfinding algorithm.

For clarity, the two relative conditions in the theorem are
\begin{equation}\label{eq:intro-relative-definitions}
 \kappa_{\mathrm{coeff,rel}}(p,z)
 =\frac{\|a\|_2\bigl(\sum_{k=0}^n|z|^{2k}\bigr)^{1/2}}
        {|z|\,|p'(z)|},\qquad
 \kappa_{\mathrm{eig,rel}}(C_p,z)
 =\Lambda(p,z)\frac{\|C_p\|_2}{|z|},
\end{equation}
where $a=(a_0,\ldots,a_n)^T$ and every polynomial coefficient is perturbable.
Theorem~\ref{thm:comp-main} proves part (i), and
Corollary~\ref{cor:comp-relative} proves part (ii).

\subsection{Proof strategy and the role of root geometry}

The proof separates local root sensitivity from the additional conditioning
of interpolation in the companion basis. To analyze local sensitivity, we
first work with systems of $d$ independent complex Gaussian polynomials
sharing a covariance kernel. That kernel defines a natural output metric.
Marked Kac--Rice shows that, after normalization by this metric, the Jacobian
at a selected root is independent of the root location and has a
determinant-weighted Gaussian law. The resulting intrinsic condition number
$\kappa_H$, defined in Section~\ref{sec:setup-kernel-metric}, satisfies
\[
 \Pstar\{\kappa_H>t\}\le
 \min\left\{1,\frac{d(d+1)}{4t^4}\right\},
\]
with the constant asymptotically exact for fixed $d$. The precise assumptions
and statements are Theorems~\ref{thm:local-universal} and~\ref{thm:local-tail}.
The induced Ginibre and Wishart distributions used here are established
random-matrix results \cite{fischmann2012induced,kumar2017smallest}.

For the Gaussian box ensemble, exact radial laws then translate this geometric
normalization back to Euclidean root displacement. They give
$\Estar\log\kappa_{\mathrm{coeff}}=-\log n+O_d(1)$ and retain the
contribution of roots outside the unit polydisk. Finally, in one variable,
writing $p(x)=(x-z)q(x)$ turns the coefficient norm into the Gram matrix of
multiplication by $x-z$. This matrix is tridiagonal. Its inverse trace, together
with the Gaussian weight at a root, gives the exact conditional overlap moment
and the bounds in Theorem~\ref{thm:intro-main}. The same argument also shows
that, for each fixed $n\ge2$, overlap moments are finite precisely below order
four; the threshold two concerns uniformity as the degree grows.

The paper follows this proof structure. Section~\ref{sec:setup} fixes the
input--output maps, metrics, and sampling law.
Section~\ref{sec:local-law} establishes the local Gaussian law;
Section~\ref{sec:box} derives the box root and coefficient scales; and
Section~\ref{sec:companion} proves the main theorem.
Section~\ref{sec:discussion} discusses its interpretation and the multivariate
questions that remain.

\subsection{Relation to previous work}\label{sec:intro-previous}

Local conditioning of polynomial systems has a substantial geometric and
probabilistic theory. Shub--Smale developed the projective framework, its
relation to ill-posedness, and distributional estimates in the
\emph{Complexity of B\'ezout's theorem} series
\cite{shub1993geometric,shub1993volumes,shub1993packing}.
B\"urgisser--Cucker \cite{burgisser2013condition} give a systematic account.
Proposition~\ref{prop:setup-condition} is a standard derivative-norm
calculation in explicitly specified input and output metrics; it is included
to fix conventions, not as a new condition-number theorem.

Beltr\'an--Pardo \cite{beltran2009smale,beltran2011fast} used average
conditioning and integration over system--root pairs to obtain average
polynomial-time homotopy algorithms. Their work also gives exact moments of
the normalized projective condition number below order four
\cite[Sections~7--8]{beltran2011fast}.
B\"urgisser--Cucker \cite{burgisser2011smale} established smoothed estimates,
and Armentano--Beltr\'an--B\"urgisser--Cucker--Shub
\cite[Theorem~2]{armentano2016condition} obtained an exact centered
inverse-Jacobian Frobenius second moment and a smoothed bound.
Section~\ref{sec:local-classical} spells out the overlap with our local law.
Lairez \cite{lairez2017deterministic} subsequently obtained a deterministic
algorithm with polynomial average complexity.

For sparse systems, Malajovich--Rojas \cite{malajovich2004high} construct
metrics from the coefficient spaces, relate conditioning to the discriminant,
and prove fourth-power tail bounds. Their normalization and worst-root
probability differ from our absolute, selected-root formulation. Toric
condition metrics and homotopy complexity are developed further in
\cite{malajovich2019toric,malajovich2023renormalization}.
Our local section presents the common-kernel Gaussian calculation in a form
suited to comparing coefficient and algebraic-representation perturbations;
the fourth-power exponent and the homogeneous moment identities already
belong to this established theory.

Toh--Trefethen \cite{toh1994pseudozeros} explain the different perturbation
geometries of polynomial roots and companion eigenvalues. Noferini--Townsend
and Graf--Townsend construct multivariate examples in which intermediate
algebraic problems amplify conditioning exponentially with the number of
variables \cite{noferini2016numerical,graf2026numerical}. Our representation
theorem concerns one variable and growing degree; the local framework extends
to systems, but the present paper does not resolve those growing-dimension
examples.

B\"urgisser--Cucker--Rocha Cardozo \cite{burgisser2015characteristic} study
the conditioning of characteristic polynomials induced by Gaussian matrices.
Their results demonstrate the importance of the induced coefficient law in
an algebraic reduction. Here the construction runs from a Gaussian polynomial
to its companion matrix, giving a different probability distribution and
perturbation problem.

\Needspace{4\baselineskip}
Armentano--Beltr\'an \cite[Theorem~1.1]{armentano2019polynomial} study the
polynomial eigenvalue problem for
$Q(z)=\sum_{j=0}^k B_jz^j$, where $B_j\in\C^{m\times m}$.
Their perturbable input is the tuple $(B_0,\ldots,B_k)$, with norm
$(\sum_j\|B_j\|_F^2)^{1/2}$, and their random model makes all entries of these
coefficient matrices independent complex Gaussians. Thus their main theorem
concerns matrix-polynomial eigenvalue conditioning; Remark~1.4 discusses
companion algorithms in this setting as well as scalar polynomial rootfinding.
The scalar specialization \cite[Section~3.1]{armentano2019polynomial} gives an
exact identity for our coefficient ensemble, using the standard Fubini--Study
output metric. In our notation,
\[
 \mu_{\mathrm{FS}}(p,z)
 =\frac{\|a\|_2\bigl(\sum_{k=0}^n|z|^{2k}\bigr)^{1/2}}
        {(1+|z|^2)|p'(z)|},\qquad \Estar\mu_{\mathrm{FS}}^2=1.
\]
This is relative to the input norm and uses a projective output metric,
explained in Section~\ref{sec:setup-kernel-metric}. It differs from our kernel
metric. The companion overlap concerns a further change of input: arbitrary
perturbations of $C_p$. Viewing $zI-C_p$ as a degree-one matrix polynomial does
not make the Gaussian coefficient-matrix theorem applicable: its leading
matrix is fixed, and $C_p$ has prescribed zeros and ones and ratios of Gaussian
coefficients. Our theorem analyzes this induced matrix ensemble. The distinction
is therefore between input spaces, metrics, and probability laws, rather than
between polynomial solving and eigenvalue problems as such.

Barbarino--Noferini study limiting empirical eigenvalue measures of random
matrix polynomials, both with a fixed identity leading coefficient
\cite{barbarino2023monic} and with random leading coefficients
\cite{barbarino2022limit}. The latter paper includes our scalar Gaussian
ensemble and its limiting uniform root distribution on the unit circle
\cite[Theorem~1.2]{barbarino2022limit}. These results describe root locations;
weak convergence of empirical measures does not determine companion overlap
moments. Their monic paper also provides bounds for the smallest singular
value of a Gaussian matrix with arbitrary deterministic mean
\cite[Theorem~A.5]{barbarino2023monic}. These are relevant tools for smoothed
analysis, although an application to root conditioning must account for
sampling at a root and for the actual perturbation structure. Both papers
treat polynomials in one variable; growing matrix size is distinct from a
growing number of variables in a polynomial system.

\section{The input--output problems}\label{sec:setup}

In this section, we specify the coefficient space and local
root map, choose the output metric, and define conditioning as a derivative
norm. We then identify the matrix perturbation problem and the probability law
used to select a root.

\subsection{Coefficient space and the local root map}
Let $\mathcal H$ be an $m$-dimensional space of complex polynomials on an open
set $U\subseteq\C^d$, with a specified Hermitian inner product and orthonormal
basis $\phi_1,\ldots,\phi_m$. Write
$\phi(z)=(\phi_1(z),\ldots,\phi_m(z))^T$. 
The input is a coefficient array $c=(c_{ik})\in\C^{d\times m}$, representing
\[
 P_c=(p_1,\ldots,p_d)^T,\qquad
 p_i(z)=\sum_{k=1}^m c_{ik}\phi_k(z).
\]
We allow perturbations of every coefficient and measure perturbations using the Frobenius norm:
\[
 \|\delta c\|_F^2=\sum_{i,k}|\delta c_{ik}|^2.
\]
Thus a change of orthonormal coefficient basis changes coordinates, but not the
input norm. A nonunitary change of polynomial basis generally changes the
problem unless its Gram matrix is included in the norm.

Write $J_c(z)=(\partial p_i/\partial z_j)_{i,j}$ for the Jacobian matrix.
A root $z$ of $P_c$ is \emph{simple} if $J_c(z)$ is invertible. We will work
with the incidence set
\[
 \mathcal V=\{(c,z)\in\C^{d\times m}\times U:
                  P_c(z)=0,\ \det J_c(z)\ne0\}.
\]
At each $(c,z)\in\mathcal V$, the implicit function theorem gives a unique local
root branch 
\begin{equation}\label{eq:setup-root-branch}
   \mathcal S_{c,z}\colon c'\mapsto z(c'), \text{ with } z(c)=z.
\end{equation}   
This local map is the input-output map under study.
Its derivative at $c$ is
\begin{equation}\label{eq:setup-root-derivative}
 D\mathcal S_{c,z}[\delta c]=-J_c(z)^{-1}E_z[\delta c],
 \qquad E_z[\delta c]=\delta c\,\phi(z).
\end{equation}
Here $E_z$ evaluates the perturbed equations at the original root:
$E_z[\delta c]=\delta P(z)$. Derivatives of the root branch and other
holomorphic maps are complex-linear.

\subsection{Choosing an output metric and defining conditioning}

To quantify a change in the output, we must choose how to measure root
displacements. Fix a smooth matrix-valued function
\[
 G:U\longrightarrow\C^{d\times d},\qquad
 G(z)=G(z)^*,\qquad v^*G(z)v>0\quad(v\ne0).
\]
Thus $G$ assigns a positive definite Hermitian matrix to every possible root
location. Identifying the tangent space at
$z\in U$ with $\C^d$, define
\[
 \langle u,v\rangle_{G,z}=u^*G(z)v,\qquad
 \|v\|_{G,z}=\sqrt{v^*G(z)v}.
\]
This field of inner products thus defines a Hermitian metric on $U$. The constant
choice $G(z)=I_d$ gives ordinary Euclidean displacement. Later we will also
use a particular nonconstant choice $G=H$, determined by the coefficient
space and defined in Section~\ref{sec:setup-kernel-metric}.

\begin{definition}
For a simple system--root pair $(c,z)$, the \emph{absolute coefficient
condition number} with this output metric is defined as 
\begin{equation}\label{eq:setup-condition-definition}
 \kappa_G(c,z)
 :=\sup_{\delta c\ne0}
    \frac{\|D\mathcal S_{c,z}[\delta c]\|_{G,z}}{\|\delta c\|_F}.
\end{equation}
\end{definition}

The squared norm of scalar polynomial evaluation at $z$ is
\[
 S(z)=\|\phi(z)\|_2^2.
\]
Thus $S(z)=0$ if and only if $z$ is a common root of all the basis polynomials
$\phi_1,\ldots,\phi_m$, or equivalently of every polynomial in $\mathcal H$.
For a positive definite Hermitian matrix $G(z)$, the notation $G(z)^{1/2}$
means its unique positive definite Hermitian square root.

The next proposition specializes the standard condition-as-derivative
framework \cite[Chapter~14]{burgisser2013condition} to our coefficient space.
We give the short proof to make the evaluation factor and the choice of
output metric explicit.

\begin{proposition}\label{prop:setup-condition}
Fix an output metric $G$ as above and a simple system--root pair $(c,z)$ with
$S(z)>0$. The condition number in \eqref{eq:setup-condition-definition} is
\begin{equation}\label{eq:setup-condition}
 \kappa_G(c,z)=\sqrt{S(z)}\,
       \|G(z)^{1/2}J_c(z)^{-1}\|_2.
\end{equation}
For $G(z)=I_d$, we write $\kappa_{\mathrm{coeff}}:=\kappa_{I_d}$, so
$\kappa_{\mathrm{coeff}}(c,z)=\sqrt{S(z)}\,\|J_c(z)^{-1}\|_2$.
For a fixed holomorphic scalar output $g:U\to\C$, with its usual absolute-value
metric, the corresponding condition number is
\begin{equation}\label{eq:setup-scalar-condition}
 \kappa_{\mathrm{coeff}}^{g}(c,z)
   =\sqrt S\,\|Dg(z)J^{-1}\|_2.
\end{equation}
\end{proposition}
\begin{proof}
Suppress $c,z$ temporarily and write $J^{-*}=(J^{-1})^*$. The evaluation map
satisfies $E_zE_z^*=S(z)I_d$. Therefore
\[
 (G^{1/2}J^{-1}E_z)(G^{1/2}J^{-1}E_z)^*
   =S\,G^{1/2}J^{-1}J^{-*}G^{1/2}.
\]
Taking operator norms proves \eqref{eq:setup-condition} and the chain rule gives
\eqref{eq:setup-scalar-condition}.
\end{proof}

The assumption $S(z)>0$ is not needed for Proposition~\ref{prop:setup-condition}:
if $S(z)=0$, then $E_z=0$, so the root derivative and the displayed condition
numbers are zero, and the formulas still hold. Positivity will be important
later when defining the kernel metric through $\log S$ and normalizing the
Jacobian by $S^{-1/2}$.

The distinction between coefficient and residual perturbations is essential.
The norm $\|J^{-1}\|_2$ measures sensitivity to perturbing the equation values
$P(z)$. Perturbing the coefficient array first evaluates a polynomial at $z$,
which contributes the factor $\sqrt S$.

For nonzero input and a nonzero Euclidean output $z$, the normwise relative condition is
\begin{equation}\label{eq:setup-relative}
 \kappa_{\mathrm{coeff,rel}}(c,z)
 =\kappa_{\mathrm{coeff}}(c,z)\frac{\|c\|_F}{\|z\|_2}.
\end{equation}
For the scalar output $g(z)\ne0$, replace the denominator by $|g(z)|$ and use
$\kappa_{\mathrm{coeff}}^g$. We do not use the phrase ``relative condition''
without specifying these input and output scales.

\subsection{The kernel metric}\label{sec:setup-kernel-metric}

We now specify the second output metric used in this paper. The coefficient
space determines the reproducing kernel
\[
 K(z,w)=\sum_{k=1}^m\phi_k(z)\overline{\phi_k(w)},\qquad S(z)=K(z,z).
\]
Throughout the general local theory we assume $S>0$ and
\begin{equation}\label{eq:setup-kernel-metric}
 H_{jk}(z)=\partial_{\bar z_j}\partial_{z_k}\log S(z),
 \qquad H(z)>0\quad(z\in U).
\end{equation}
The derivatives are taken with $z$ and $\bar z$ as independent variables
(Wirtinger derivatives). The order of the indices is chosen so that a tangent
vector $v$ has squared length $v^*H(z)v$. Taking $G=H$ in the preceding
subsection defines the kernel output metric and its condition number $\kappa_H$.

There is a geometric interpretation of this formula. The notation $[\phi(z)]$
means the complex line spanned by the nonzero vector $\phi(z)$, a point of
complex projective space. The Fubini--Study metric is the standard Hermitian
metric on that space of lines. Its \emph{pullback} by $z\mapsto[\phi(z)]$
measures a displacement at $z$ by the displacement it induces in this line;
the resulting matrix is $H(z)$. Positivity means that the line changes to
first order in every nonzero tangent direction. We use the normalization
whose expression on $\mathbb P^1$ is $|\dd z|^2/(1+|z|^2)^2$.
The explicit formula \eqref{eq:setup-kernel-metric} suffices to use the metric
without this projective interpretation.

Constructing the root metric from the coefficient inner product is also
central to the sparse-system framework of Malajovich--Rojas
\cite[Sections~1--2]{malajovich2004high}; input normalizations must be
matched before comparing condition numbers.

Define the normalized Jacobian
\begin{equation}\label{eq:setup-normalized}
 A=S^{-1/2}JH^{-1/2}.
\end{equation}
Proposition~\ref{prop:setup-condition} gives the exact identity
\begin{equation}\label{eq:setup-kernel-condition}
 \kappa_H(c,z)=\|A^{-1}\|_2.
\end{equation}
This is a change of output metric, with the coefficient norm held fixed.
The Euclidean and kernel-metric condition numbers answer different questions;
Section~\ref{sec:box} relates them in the box model.

\begin{example}[Two natural metrics in one variable]\label{ex:setup-metrics}
For the monomial coefficient norm on polynomials of degree at most $n$,
$S(z)=\sum_{k=0}^n|z|^{2k}$. With $r=|z|^2$,
\[
 H(z)=h_n(r)=\left(r\frac{k_n'(r)}{k_n(r)}\right)',
 \qquad k_n(r)=\sum_{k=0}^n r^k.
\]
For the different coefficient norm making
$\sqrt{\binom nk}z^k$ orthonormal, $S(z)=(1+|z|^2)^n$ and
$H(z)=n/(1+|z|^2)^2$. The latter is the complex Kostlan model after Gaussian
coefficients are introduced. The universal local law below applies to both,
but their kernel metrics and their coefficient norms differ.
\end{example}

For comparison with the classical projective theory, consider $d$ homogeneous
equations of a common degree $n$ in $d+1$ variables with the
Bombieri--Weyl inner product. This makes the monomials
$\sqrt{n!/\alpha!}\,X^\alpha$, $|\alpha|=n$, orthonormal, where
$\alpha!=\alpha_0!\cdots\alpha_d!$.
In the affine chart $X_0=1$, their kernel satisfies
$S(z)=(1+\|z\|_2^2)^n$, so $H=nG_{\mathrm{FS}}$, with
$G_{\mathrm{FS}}$ the standard projective metric in this chart.
The classical degree-normalized projective condition number is therefore
\begin{equation}\label{eq:setup-classical-mu}
 \mu(c,[1:z])=\|c\|_F\,\kappa_H(c,z).
\end{equation}
Indeed, at $z=0$ the normalized projective Jacobian is $J/\sqrt n$;
unitary invariance gives the identity in every chart
\cite{shub1993geometric,burgisser2013condition}.
Thus $\mu$ includes both an input-norm factor and a degree normalization;
the coefficient derivative measured in $G_{\mathrm{FS}}$ has norm
$\kappa_H/\sqrt n$. Neither should be identified with the affine relative
condition in \eqref{eq:setup-relative}.

\subsection{The matrix eigenvalue problem and its perturbations}

For an eigenvalue $\lambda$ of algebraic multiplicity one of
$M\in\C^{N\times N}$, choose nonzero
vectors $r,w$ such that $Mr=\lambda r$ and $w^*M=\lambda w^*$.
The local eigenvalue map $\mathcal E_{M,\lambda}$ sends a nearby matrix to
its unique eigenvalue near $\lambda$. Its derivative at $M$ is
\[
 D\mathcal E_{M,\lambda}[\delta M]=\frac{w^*\delta M r}{w^*r}.
\]
With arbitrary matrix perturbations measured in spectral norm and scalar
output measured in absolute value, its condition is
\begin{equation}\label{eq:setup-overlap}
 \kappa_{\mathrm{eig}}(M,\lambda)
 =\Lambda(M,\lambda):=\frac{\|w\|_2\|r\|_2}{|w^*r|}\ge1.
\end{equation}
The derivative norm is at most the displayed overlap, and a rank-one
perturbation attains it; $\Lambda\ge1$ follows from Cauchy--Schwarz.
The same absolute condition is obtained using Frobenius norm,
but its relative version would then use $\|M\|_F$. Here we consistently use
\begin{equation}\label{eq:setup-eigen-relative}
 \kappa_{\mathrm{eig,rel}}(M,\lambda)
 =\Lambda(M,\lambda)\frac{\|M\|_2}{|\lambda|}
 \quad(\lambda\ne0).
\end{equation}

Suppose a differentiable algebraic construction $\mathcal D:c\mapsto M(c)$
represents $g(z)$ as an eigenvalue. On compatible local branches,
\begin{equation}\label{eq:setup-composition}
 \mathcal E\circ\mathcal D=g\circ\mathcal S,
 \qquad D\mathcal E\,D\mathcal D=Dg\,D\mathcal S.
\end{equation}
Consequently coefficient-induced matrix perturbations give exactly the
original derivative. The comparison in this paper instead allows arbitrary
perturbations in the intermediate matrix space. It describes the sensitivity
of that larger problem. Errors in constructing $M$, implementing an eigensolver,
and recovering other coordinates require additional analysis.

One important construction uses the quotient algebra of the roots. Suppose
the ideal $I=\langle p_1,\ldots,p_d\rangle$ has exactly $N$ distinct complex
zeros $z_1,\ldots,z_N$ and consists of all polynomials vanishing at these
points (equivalently, $I$ is radical). In
$\mathcal A=\C[x_1,\ldots,x_d]/I$, two polynomials represent the same element
precisely when their values at all roots agree. Evaluation identifies
$\mathcal A$ with the $N$-dimensional space of functions on this root set.
Choose a basis $\B=(b_1,\ldots,b_N)$ of polynomial representatives.
For a polynomial $g$, multiplication by its class is a linear map on
$\mathcal A$; denote its matrix, acting on coefficient columns, by $M_g$.
The evaluation matrix $V=(b_k(z_j))_{k,j}$ is invertible, and evaluation gives
\[
 M_g=V^{-T}\diag(g(z_1),\ldots,g(z_N))V^T.
\]
If $g$ \emph{separates} the roots, meaning that $g(z_1),\ldots,g(z_N)$ are
pairwise distinct, its overlap at $z_j$ is therefore
\begin{equation}\label{eq:setup-interpolation}
 \Lambda_j=\|V^{-T}e_j\|_2\|Ve_j\|_2,
 \qquad
 \Lambda_j^2=
 \frac{\|Ve_j\|_2^2}
 {\dist(Ve_j,\spanop\{Ve_k:k\ne j\})^2}.
\end{equation}
The last identity follows by orthogonally projecting one evaluation column
away from the others. Thus matrix conditioning becomes a question about
interpolation in a specified basis norm. This deterministic interpretation
also underlies the normal-form analysis in \cite{graf2026numerical}.

\subsection{Gaussian input and a uniformly selected root}\label{sec:setup-probability}

We now make the input random: all $dm$ coefficient coordinates are independent
$\CN(0,1)$ variables, with density $\pi^{-1}e^{-|a|^2}$ on $\C$.
Equivalently, the real and imaginary parts have variance $1/2$.
The components of $P$ are independent and share the kernel $K$.

Assume that almost surely $P$ has exactly $N\ge1$ zeros in $U$, all simple,
where $N$ is deterministic. First sample $P$ with this Gaussian law, then select
one of its roots uniformly. For a nonnegative measurable function $F(P,z)$,
called a \emph{mark}, define
\begin{equation}\label{eq:setup-selected-root}
 \Estar F(P,Z)=\frac1N\E\sum_{P(z)=0,\ z\in U}F(P,z).
\end{equation}
This determines a probability law $\Pstar$ on system--root pairs. In particular,
the marginal law of $P$ is its original Gaussian law. This last fact uses the
deterministic root count. No independence between the roots of one system is
assumed.

Sampling a root this way is different from conditioning the Gaussian input on
$P(z)=0$ at a prescribed point. The former also weights the conditional law
by $|\det J(z)|^2$, as Section~\ref{sec:local-law} proves. We use conditional
expectations given $Z=z$ in the regular conditional sense supplied by that
formula, rather than treating a zero-probability event as an ordinary event.

For example, for an overlap threshold $t$,
\[
 \Pstar(\Lambda>t)=\frac1N\E\#\{j:\Lambda_j>t\}.
\]
A tail bound therefore controls the expected fraction of problematic roots.
It does not give the same bound for the worst root of a system. Likewise
$\Estar\log X$ is a logarithmic average, not a median or a high-probability
bound for $X$.

\section{The normalized Jacobian at a selected root}
\label{sec:local-law}

We use the common-kernel Gaussian model and the selected-root law of
Section~\ref{sec:setup-probability}, with $S>0$, $H>0$, and exactly $N\ge1$
simple roots in $U$ almost surely. The same argument applies to a finite
collection of holomorphic functions under these assumptions. Write $dV$
for Lebesgue measure on $\C^d\simeq\mathbb R^{2d}$.

\subsection{Gaussian regression and the kernel metric}

The order of the indices in \eqref{eq:setup-kernel-metric} fixes our complex covariance
convention.  A derivative $Dp(z)$ is a row, whereas a root displacement is a
column.  With $D\phi(z)$ the matrix whose $k$th column is
$\partial_{z_k}\phi(z)$, direct differentiation gives
\begin{equation}\label{eq:local-metric-projection}
  H(z)=\frac1{S(z)}D\phi(z)^*
  \left(I-\frac{\phi(z)\phi(z)^*}{S(z)}\right)D\phi(z).
\end{equation}
Thus $\delta z^*H(z)\delta z$ is the squared norm of the derivative of
$z\mapsto[\phi(z)]$ in the Fubini--Study metric, with the normalization specified
by \eqref{eq:local-metric-projection}.  Positivity of $H$ says that this
projective evaluation map is an immersion.

\begin{lemma}[Conditional derivative law]\label{lem:local-regression}
  At any fixed $z\in U$,
  \[
    \E\bigl[Dp_i(z)^*Dp_i(z)\mid p_i(z)=0\bigr]=S(z)H(z).
  \]
  Consequently, conditionally on $P_c(z)=0$,
  \begin{equation}\label{eq:local-regression}
    J_c(z)\ \stackrel{\mathrm d}{=}\ \sqrt{S(z)}\,G_0 H(z)^{1/2},
  \end{equation}
  where $G_0$ has independent $\CN(0,1)$ entries and the square root is the
  positive Hermitian square root.
\end{lemma}

\begin{proof}
  For one component, Gaussian regression applied to $(p,Dp)$ gives
  \[
    \E\bigl[\overline{\partial_jp}\,\partial_kp\mid p=0\bigr]
    =\partial_{\bar z_j}\partial_{z_k}S
     -\frac{(\partial_{\bar z_j}S)(\partial_{z_k}S)}{S}
    =S H_{jk}.
  \]
  The conditional mean and the conditional pseudo-covariance are zero.
  The components remain independent after conditioning on their respective
  values.  A standard Gaussian row $g$ satisfies
  $\E[(gH^{1/2})^*(gH^{1/2})]=H$, which proves
  \eqref{eq:local-regression} with our row convention.
\end{proof}

Recall that $A=S^{-1/2}JH^{-1/2}$ and $\kappa_H=\|A^{-1}\|_2$.
Although $A$ is a standard Gaussian matrix under conditioning at a fixed point,
its law at a selected root has an additional determinant weight.

\subsection{The marked Kac--Rice law}

Integration over system--root pairs by coarea is standard in polynomial
conditioning; see \cite[Sections~5--6]{beltran2011fast} and
\cite[Remark~5]{armentano2016condition}. The determinant weight below is its
complex Kac--Rice form.

\begin{lemma}[Marked Kac--Rice]\label{lem:local-rice}
  For every nonnegative Borel function $F$ of the full coefficient array and
  the point,
  \begin{equation}\label{eq:local-rice}
    \E\sum_{z\in U:\,P_c(z)=0}F(c,z)
    =\int_U\frac{
       \E\bigl[F(c,z)|\det J_c(z)|^2\mid P_c(z)=0\bigr]
      }{(\pi S(z))^d}\,dV(z).
  \end{equation}
  The equality is valid in $[0,\infty]$.  It also holds for signed or complex
  marks whenever either side with $|F|$ is finite.
\end{lemma}

\begin{proof}
  Regard $P_c$ as a smooth Gaussian map between real spaces of dimension $2d$.
  Its value at $z$ has density $(\pi S(z))^{-d}$ at zero, and holomorphy gives
  $\det D_{\mathbb R}P_c(z)=|\det J_c(z)|^2$.  The value distribution is
  nondegenerate since $S>0$, and the assumed simplicity excludes singular
  zeros almost surely.  The Rice formula therefore applies on compact
  subsets of $U$.  For a bounded continuous mark, use the full finite
  coefficient array as the auxiliary Gaussian field in the weighted Rice
  formula \cite[Theorems~6.2 and~6.4]{azais2009level}.
  A monotone-class argument extends the identity to bounded Borel marks.
  Truncation of the mark and exhaustion of $U$ by compact sets then give
  \eqref{eq:local-rice} by monotone convergence.  Finally apply the
  nonnegative identity to $|F|$ to justify the integrable case.
\end{proof}

\begin{theorem}[Universal local law]\label{thm:local-universal}
  Under $\Pstar$, the root location $Z$ and the normalized Jacobian $A(c,Z)$
  are independent.  Their densities are respectively
  \begin{equation}\label{eq:local-root-density}
    \frac{d!}{N\pi^d}\det H(z),\qquad z\in U,
  \end{equation}
  with respect to $dV(z)$, and
  \begin{equation}\label{eq:local-induced-density}
    q_d(A)=\frac{1}{d!\,\pi^{d^2}}
      |\det A|^2\exp\bigl(-\|A\|_F^2\bigr),
      \qquad A\in\C^{d\times d},
  \end{equation}
  with respect to Lebesgue measure on the matrix entries.  Moreover,
  \begin{equation}\label{eq:local-Wishart}
    A^*A\ \stackrel{\mathrm d}{=}\ B^*B,
    \qquad B\in\C^{(d+1)\times d}
  \end{equation}
  for $B$ with independent $\CN(0,1)$ entries.
\end{theorem}

\begin{proof}
  Apply Lemma~\ref{lem:local-rice} with
  $F(c,z)=f(z)g(A(c,z))$, where $f,g\geq0$.  Lemma~\ref{lem:local-regression}
  and
  \[
    |\det J_c(z)|^2=S(z)^d\det H(z)\,|\det A(c,z)|^2
  \]
  yield
  \begin{equation}\label{eq:local-factorization}
    \Estar[f(Z)g(A)]
    =\frac{1}{N\pi^d}\int_U f(z)\det H(z)\,dV(z)
      \E\bigl[g(G_0)|\det G_0|^2\bigr].
  \end{equation}
  Expanding the determinant over permutations gives
  $\E|\det G_0|^2=d!$: only equal pairs of permutations have nonzero
  expectation.  With $f=g=1$, \eqref{eq:local-factorization} shows that
  \eqref{eq:local-root-density} integrates to one.  The factorization then
  proves independence and \eqref{eq:local-induced-density}.

  A complex Wishart matrix $B^*B$, with $B$ of size $m\times d$, has density
  proportional to $(\det W)^{m-d}e^{-\operatorname{tr}W}$ on positive
  Hermitian matrices.  The square Gaussian Gram matrix $G_0^*G_0$ has exponent
  zero.  Multiplication of its density by $|\det G_0|^2=\det(G_0^*G_0)$ therefore
  gives precisely the case $m=d+1$, proving \eqref{eq:local-Wishart}.
\end{proof}

The matrix law \eqref{eq:local-induced-density} is the complex induced Ginibre
ensemble with parameter one \cite{fischmann2012induced}. The zero-density
formula belongs to the established covariance-kernel approach to Gaussian
zeros \cite{bleher2000universality,malajovich2004high}. These are standard
ingredients. Their
role here is to identify exactly the distribution of the coefficient-to-root
derivative in the kernel metric, separating it from the distribution of root
locations and from any algebraic representation.

\subsection{A sharp tail for local conditioning}

\begin{theorem}[Tail and moments]\label{thm:local-tail}
  Put $C_d=d(d+1)/4$.  For every $t>0$,
  \begin{equation}\label{eq:local-tail}
    \Pstar\{\kappa_H>t\}\leq\min\{1,C_dt^{-4}\},
    \qquad
    \lim_{t\to\infty}t^4\Pstar\{\kappa_H>t\}=C_d.
  \end{equation}
  In particular, for $q>0$,
  \begin{equation}\label{eq:local-moment-range}
    \Estar\kappa_H^q<\infty\quad\Longleftrightarrow\quad q<4,
    \qquad
    \Estar\kappa_H^q\leq\frac{4}{4-q}C_d^{q/4}
    \quad(0<q<4).
  \end{equation}
  The inverse Frobenius moment is exactly
  \begin{equation}\label{eq:local-frobenius}
    \Estar\|A^{-1}\|_F^2=d.
  \end{equation}
  Also $\Estar|\log\kappa_H|<\infty$.
\end{theorem}

\begin{proof}
  Let $X=\sigma_{\min}(A)^2$.  By \eqref{eq:local-Wishart}, this is the smallest
  eigenvalue of a complex Wishart matrix with rectangularity one.  Its density
  is \cite[Section~2.2, equation~(15)]{kumar2017smallest}
  \begin{equation}\label{eq:local-smallest-density}
    f_d(x)=x e^{-dx}L_{d-1}^{(2)}(-x)
    =x e^{-dx}\sum_{k=0}^{d-1}
        \binom{d+1}{k+2}\frac{x^k}{k!},\qquad x>0.
  \end{equation}
  For $0\leq k\leq d-1$,
  \[
    \frac{\binom{d+1}{k+2}}{\binom{d+1}{2}}
      =\frac{2(d-1)(d-2)\cdots(d-k)}{(k+2)!}
      \leq d^k,
  \]
  where the empty product is one.  Hence
  $f_d(x)\leq\binom{d+1}{2}x$.  Since
  $\kappa_H=X^{-1/2}$, integration over $0<x<t^{-2}$ proves the bound in
  \eqref{eq:local-tail}.  The expansion
  $f_d(x)=\binom{d+1}{2}x+O_d(x^2)$ at zero proves its sharp asymptotic.
  Integrating the tail gives both the moment threshold and, on splitting at
  $t=C_d^{1/4}$, the bound in \eqref{eq:local-moment-range}.

  There is an elementary cofactor proof of \eqref{eq:local-frobenius}:
  \[
    \Estar\|A^{-1}\|_F^2
      =\frac{\E\bigl[|\det G_0|^2\|G_0^{-1}\|_F^2\bigr]}{d!}
      =\frac{\E\|\operatorname{adj}G_0\|_F^2}{d!}
      =\frac{d^2(d-1)!}{d!}=d.
  \]
  Each cofactor is the determinant of a $(d-1)\times(d-1)$ standard complex
  Gaussian matrix and has second moment $(d-1)!$; at $d=1$ the empty
  determinant is one.  Finally, the positive logarithmic part follows from
  any positive moment below four.  For the negative part, use
  $\log^+(1/\kappa_H)=\log^+\sigma_{\min}(A)
   \leq\log^+\|A\|_F$ and the Gaussian decay in
  \eqref{eq:local-induced-density}.
\end{proof}

Two consequences clarify the scope of the theorem.  First,
\begin{equation}\label{eq:local-dimension-tail}
  \Pstar\{\kappa_H>\sqrt d\,t\}
  \leq\min\{1,(2t^4)^{-1}\}.
\end{equation}
Thus the geometric normalization gives a tail uniform in the coefficient
space, polynomial degrees, and, after the scale $\sqrt d$, the dimension.
Second, \eqref{eq:setup-selected-root} makes the same probability an expected
fraction of roots:
\[
  \frac1N\E\#\{z\in U:P_c(z)=0,\ \kappa_H(c,z)>t\}
  =\Pstar\{\kappa_H>t\}.
\]
This statement does not assume independence between different roots.

\subsection{Comparison with classical projective moments}
\label{sec:local-classical}

In the common-degree Bombieri--Weyl model of
\eqref{eq:setup-classical-mu}, the classical Frobenius variant is
$\mu_F(c,[1:z])=\|c\|_F\|A^{-1}\|_F$.
Consequently \eqref{eq:local-frobenius} recovers the centered identity
\begin{equation}\label{eq:local-known-frobenius}
 \Estar\frac{\mu_F(c,[1:Z])^2}{\|c\|_F^2}=d
\end{equation}
of Armentano--Beltr\'an--B\"urgisser--Cucker--Shub
\cite[Theorem~2]{armentano2016condition}. Their theorem also allows unequal
degrees and gives the bound $d/\sigma^2$ for Gaussian perturbations of an
arbitrary deterministic homogeneous system, with equality at zero center.
Thus this homogeneous second-moment identity is already known.

Beltr\'an--Pardo \cite[Sections~7--8]{beltran2011fast} give exact operator
moments of $\mu$ for orders $0<q<4$ under spherical Bombieri--Weyl input.
Those results use the same rectangular Gaussian singular-value mechanism.
Malajovich--Rojas \cite[Theorem~1]{malajovich2004high} already obtain a
fourth-power tail for the event that some sparse-system root has large
condition number, with their projective input normalization.
Theorem~\ref{thm:local-tail} gives the explicit distribution and sharp
constant for the absolute kernel-metric condition at a selected root under
our common-kernel assumptions. Its moment threshold is consistent with these
earlier results; it does not establish a new exponent for polynomial systems.

\begin{example}[One complex equation]\label{ex:local-one-variable}
  When $d=1$, the variable $|A|^2$ has density $x e^{-x}$ on $x>0$.
  For every nondegenerate finite-dimensional Gaussian polynomial ensemble
  covered by the hypotheses,
  \[
    \Pstar\{\kappa_H>t\}
      =1-e^{-t^{-2}}(1+t^{-2}),\qquad
    \Estar\kappa_H^2=1.
  \]
  The tail is asymptotic to $1/(2t^4)$, independently of the degree.
\end{example}

There is also a short consistency check for the known projective identity
$\Estar\mu_{\mathrm{FS}}^2=1$ in the scalar monomial model
\cite{armentano2019polynomial}. Kac--Rice cancels the derivative denominator,
while the conditional coefficient vector given $p(z)=0$ is a standard
Gaussian on an $n$-dimensional hyperplane. Hence
\[
 \Estar\mu_{\mathrm{FS}}^2
 =\frac1{n\pi}\int_{\C}
    \frac{\E[\|a\|_2^2\mid p(z)=0]}{(1+|z|^2)^2}\,dV(z)
 =\frac1\pi\int_{\C}\frac{dV(z)}{(1+|z|^2)^2}=1.
\]
This identity and Example~\ref{ex:local-one-variable} concern different
condition numbers, despite their equal second moments.

\section{The box ensemble: root location and coefficient scale}\label{sec:box}

Fix $d,n\ge1$ and let
\begin{equation}\label{eq:box-model}
 p_i(z)=\sum_{\alpha\in\{0,\ldots,n\}^d}c_{i,\alpha}z^\alpha,
 \qquad c_{i,\alpha}\sim\CN(0,1)\quad\hbox{independently}.
\end{equation}
The monomials are orthonormal in the coefficient norm. For generic coefficients,
the system has
\begin{equation}\label{eq:box-count}
 N=d!n^d
\end{equation}
simple affine roots, all in $(\C^*)^d$. This is the mixed-volume count for
$d$ identical boxes $[0,n]^d$ \cite{bernstein1975roots,cox2005using}.
The exceptional coefficient sets are proper algebraic sets and hence have
Gaussian measure zero. Equivalently, the multihomogeneous compactification
on $(\mathbb P^1)^d$ has intersection number $d!n^d$, with no boundary
intersection for generic sections. This statement concerns that
compactification; ordinary homogenization to total degree $dn$ can introduce
common zeros at projective infinity.

Put
\[
 k_n(r)=\sum_{j=0}^n r^j,\qquad
 h_n(r)=\left(r\frac{k_n'(r)}{k_n(r)}\right)'.
\]
The covariance kernel and its metric are
\begin{equation}\label{eq:box-kernel}
 K(z,w)=\prod_{\ell=1}^d k_n(z_\ell\bar w_\ell),\qquad
 H(z)=\diag(h_n(|z_1|^2),\ldots,h_n(|z_d|^2)).
\end{equation}
In particular $H>0$. A convenient interpretation is obtained by giving
$J_r\in\{0,\ldots,n\}$ probabilities
$\Prob(J_r=j)=r^j/k_n(r)$. Differentiation gives
\begin{equation}\label{eq:box-variance}
 r h_n(r)=\Var(J_r),\qquad r>0.
\end{equation}
The continuous endpoint values are $h_n(0)=1$ and
$h_n(1)=n(n+2)/12$.

\subsection{An exact selected-root law}

\begin{proposition}[Independent radial coordinates]\label{prop:box-root-law}
Under $\Pstar$, the coordinates of $Z$ are independent, their angles are
independent and uniform, and $R_\ell=|Z_\ell|^2$ has distribution
\begin{equation}\label{eq:box-cdf}
 F_n(r)=\Pstar(R_\ell\le r)=\frac{r k_n'(r)}{n k_n(r)},\qquad
 F_n'(r)=\frac{h_n(r)}n\quad(r\ge0).
\end{equation}
Moreover, $R_\ell$ and $R_\ell^{-1}$ have the same law.
\end{proposition}
\begin{proof}
Theorem~\ref{thm:local-universal} and \eqref{eq:box-count} give joint density
$\prod_\ell h_n(|z_\ell|^2)/(n\pi)$ with respect to complex Lebesgue measure.
In polar coordinates $\dd V(z_\ell)=\tfrac12\dd r_\ell\dd\theta_\ell$,
which proves the density statement. Integration gives \eqref{eq:box-cdf},
with limits zero and one since $r k_n'/k_n$ tends to zero and $n$ at the
two endpoints. The identity $k_n(r)=r^n k_n(r^{-1})$ implies
$F_n(r)=1-F_n(r^{-1})$ for $r>0$, proving inversion symmetry.
\end{proof}

The independence in this proposition is between coordinates of one selected
root. Different roots of the same polynomial system remain dependent.

\begin{lemma}[Depth and the outer-root moment]\label{lem:box-depth}
For any one coordinate, put $s=n\log R$. Then
\begin{equation}\label{eq:box-depth}
 \Pstar(|s|>t)\le\min\{1,2/t\}\quad(t>0),\qquad
 \Estar\log^+R=\frac{\log(n+1)}n.
\end{equation}
Consequently $\sup_n\Estar\log(1+|s|)<\infty$.
\end{lemma}
\begin{proof}
For $0<r<1$, the mean of the truncated geometric law in
\eqref{eq:box-variance} is at most the untruncated mean $r/(1-r)$.
Thus
\[
 F_n(e^{-t/n})\le\frac1{n(e^{t/n}-1)}\le\frac1t.
\]
Inversion symmetry gives the two-sided tail. Also
\begin{align*}
 \Estar\log^+R
 &=\int_1^\infty\frac{1-F_n(r)}r\dd r\\
 &=\frac1n\left[\log\frac{r^n}{k_n(r)}\right]_{r=1}^{r=\infty}
 =\frac{\log(n+1)}n.
\end{align*}
Finally integrate the depth tail against $(1+t)^{-1}\dd t$.
\end{proof}

\subsection{Uniform metric estimates}

The next elementary estimate is used both for coefficient conditioning and
for the companion overlap. It replaces a boundary-layer asymptotic by a bound
valid at every radius and degree.

\begin{lemma}[The one-variable kernel metric]\label{lem:box-metric-bound}
There are absolute constants $0<c<C<\infty$ such that, for $n\ge1$ and
$0\le r\le1$,
\begin{equation}\label{eq:box-metric-bound}
 \frac{c}{(1-r+1/n)^2}\le h_n(r)
 \le\frac{C}{(1-r+1/n)^2}.
\end{equation}
For $r>0$,
\begin{equation}\label{eq:box-metric-reversal}
 h_n(r)=r^{-2}h_n(r^{-1}).
\end{equation}
\end{lemma}
\begin{proof}
For $0\le r\le1/4$, the pair of indices $0,1$ in the variance formula gives
$h_n(r)\ge k_n(r)^{-2}\ge(1-r)^2\ge9/16$. Direct differentiation of
$k_n(r)=(1-r^{n+1})/(1-r)$ gives, for $0<r<1$,
\begin{equation}\label{eq:box-explicit-h}
 h_n(r)=\frac1{(1-r)^2}
       -\frac{(n+1)^2r^n}{(1-r^{n+1})^2}
 \le\frac1{(1-r)^2}.
\end{equation}
Thus $h_n$ is bounded above and below by absolute positive constants on
$[0,1/4]$, which proves the assertion there.

For $1/4\le r<1$, set
\[
 m=\min\{n,\lfloor(1-r)^{-1}\rfloor\}.
\]
Then $k_n(r)\le m+1$. Also $\log r\ge-(1-r)/r\ge-4(1-r)$, so
$r^j\ge e^{-4}$ for $0\le j\le m$. Each of these indices has probability
at least $e^{-4}/(m+1)$. Using an independent copy $J_r'$,
\[
 \Var(J_r)=\frac12\E(J_r-J_r')^2
 \ge e^{-8}\Var(\operatorname{Unif}\{0,\ldots,m\})
 =e^{-8}\frac{m(m+2)}{12}.
\]
Since $m\ge\tfrac12\min\{n,(1-r)^{-1}\}$, division by $r\le1$
proves the lower bound in \eqref{eq:box-metric-bound}.
For the upper bound, $\Var(J_r)\le n^2/4$ and $r\ge1/4$ imply
$h_n(r)\le n^2$. Combining this with \eqref{eq:box-explicit-h} gives
\[
 h_n(r)\le\min\{n^2,(1-r)^{-2}\}
 \le\frac4{(1-r+1/n)^2}.
\]
At $r=1$, use $h_n(1)=n(n+2)/12$. Finally, differentiating
$F_n(r)=1-F_n(r^{-1})$ proves \eqref{eq:box-metric-reversal}.
\end{proof}

\begin{corollary}[Logarithmic metric scale]\label{cor:box-metric-log}
For $R$ with distribution \eqref{eq:box-cdf},
\begin{equation}\label{eq:box-metric-log}
 \Estar\left|\log\frac{h_n(R)}{n^2}\right|\le C
\end{equation}
with an absolute constant. In particular $\Estar\log h_n(R)=2\log n+O(1)$.
\end{corollary}
\begin{proof}
For $0<r\le1$, write $s=n\log r\le0$. Since
$n(1-r)\le|s|$, Lemma~\ref{lem:box-metric-bound} gives
\[
 \frac{c}{(1+|s|)^2}\le\frac{h_n(r)}{n^2}\le C.
\]
For $r>1$, apply reversal. At every $r>0$ this yields
\[
 \left|\log\frac{h_n(r)}{n^2}\right|
 \le C+2\log(1+|n\log r|)+2\log^+r.
\]
Take expectations and apply Lemma~\ref{lem:box-depth}, noting that
$\log(n+1)/n\le\log2$.
\end{proof}

\subsection{Coefficient conditioning and representation overlap}

\begin{proposition}[The coefficient scale]\label{prop:box-coefficient-scale}
For fixed $d$ in the box model,
\begin{equation}\label{eq:box-coefficient-scale}
 \Estar\log\kappa_{\mathrm{coeff}}(P,Z)=-\log n+O_d(1).
\end{equation}
The same formula holds for the matched scalar coordinate condition
$\kappa_{\mathrm{coeff}}^{x_1}$ in \eqref{eq:setup-scalar-condition}.
All logarithms here are absolutely integrable for each $n,d$.
\end{proposition}
\begin{proof}
Put $H=n^2D$, with diagonal entries $h_n(R_\ell)/n^2$. From
$J=\sqrt S A H^{1/2}$ we obtain
\[
 \kappa_{\mathrm{coeff}}=n^{-1}\|D^{-1/2}A^{-1}\|_2.
\]
The logarithm of $\|D^{-1/2}A^{-1}\|_2$ differs from that of
$\|A^{-1}\|_2$ by at most
$\tfrac12\max_\ell|\log D_{\ell\ell}|$.
Theorem~\ref{thm:local-tail} and Corollary~\ref{cor:box-metric-log}
give \eqref{eq:box-coefficient-scale} and absolute integrability.
For the scalar output there is the more precise factorization
\begin{equation}\label{eq:box-coordinate-factorization}
 \kappa_{\mathrm{coeff}}^{x_1}
 =h_n(R_1)^{-1/2}\|e_1^T A^{-1}\|_2.
\end{equation}
The root location and $A$ are independent. The row norm has an integrable
logarithm because it is at most $\|A^{-1}\|_2$ and at least
$\|A\|_2^{-1}$. Apply Corollary~\ref{cor:box-metric-log} again.
\end{proof}

Suppose now that a fixed monomial set $\B$ is a quotient basis on a generic
coefficient chart. Use multiplication by $x_1$, whose root values are distinct
almost surely in the box model. One way to check the last assertion is to
consider pairs $z\ne w$ with $z_1=w_1$: evaluation at these points gives two
independent linear conditions on each polynomial, since constants and all
coordinates belong to the coefficient space. The pair incidence set has
dimension at most one less than coefficient space. Its projection therefore
has Gaussian measure zero.

Define the extended nonnegative expectation
$\Theta_{n,d}^{\B}=\Estar\log\Lambda_{\B}$, where the overlap uses the
Euclidean norm on the chosen quotient coordinates. Equations
\eqref{eq:setup-overlap} and \eqref{eq:box-coefficient-scale} imply
\begin{equation}\label{eq:box-decomposition}
 \Estar\log\frac{\kappa_{\mathrm{eig}}(M_{x_1},Z_1)}
                     {\kappa_{\mathrm{coeff}}^{x_1}(P,Z)}
 =\log n+\Theta_{n,d}^{\B}+O_d(1).
\end{equation}
This identity also holds with the full-root condition in the denominator,
but \eqref{eq:box-decomposition} compares the same scalar output.
If $\Theta_{n,d}^{\B}=+\infty$, interpret the equality in the extended sense;
the other logarithmic terms are integrable. No bound on multivariate overlap
is asserted here. Section~\ref{sec:companion} controls it for the univariate
companion basis.

For completeness, the evaluation factor itself satisfies
\begin{equation}\label{eq:box-evaluation-scale}
 \Estar\log k_n(R)=2\log n+O(1),\qquad
 \Estar\log\sqrt{S(Z)}=d\log n+O_d(1).
\end{equation}
Indeed reversal gives the exact identity
\[
 \log k_n(R)=\log k_n(e^{-|s|/n})+s_+.
\]
Here $\Estar s_+=\log(n+1)$, while
\[
 \frac{n}{1+|s|}\le k_n(e^{-|s|/n})\le n+1.
\]
The lower bound follows by comparing the sum with
$\int_0^n e^{-|s|x/n}\dd x$ and using $1-e^{-u}\ge u/(1+u)$.
Lemma~\ref{lem:box-depth} proves \eqref{eq:box-evaluation-scale}.
Replacing coefficient conditioning by residual conditioning would therefore
change the leading term in \eqref{eq:box-decomposition} from $\log n$ to
$(d+1)\log n$. The outer roots contribute to this difference, even though
most roots lie close to the unit polytorus.

\section{Conditioning of companion eigenvalues}
\label{sec:companion}

We now keep the scalar output fixed and change its representation.  Let
\[
  p(x)=\sum_{j=0}^{n}a_jx^j,
  \qquad a_0,\ldots,a_n\ \text{independent with law }\CN(0,1),
\]
and let $C_p$ be multiplication by $x$ on $\C[x]/(p)$ in the ordered
basis $1,x,\ldots,x^{n-1}$.  Thus $C_p$ has ones on its subdiagonal and
last column $-(a_0,\ldots,a_{n-1})^{T}/a_n$.  Almost surely $a_n\ne0$
and all roots are simple.  For a root $z$, put
\[
  v(z)=(1,z,\ldots,z^{n-1})^T,
  \qquad
  \ell_z(x)=\frac{p(x)}{(x-z)p'(z)}.
\]
The coefficient vector of $\ell_z$ is a right eigenvector of $C_p$,
whereas $v(z)^T$ is a left eigenvector, and $\ell_z(z)=1$.  Consequently
the absolute eigenvalue condition number for spectral-norm matrix
perturbations is
\begin{equation}\label{eq:comp-overlap}
  \Lambda(p,z)=\|v(z)\|_2\,\|\ell_z\|_2.
\end{equation}
Here the second norm is the Euclidean norm of the coefficient vector in
the displayed monomial basis.  In particular $\Lambda\ge1$.  This basis
choice is part of the theorem: the overlap is not invariant under a
general change of polynomial basis.

The next theorem proves part (i) of Theorem~\ref{thm:intro-main}.
The comparison in part (ii) follows in Corollary~\ref{cor:comp-relative}.

\begin{theorem}[Uniform companion-overlap tails]\label{thm:comp-main}
There is an absolute constant $C$ such that, for every $n\ge1$ and $t>0$,
\begin{equation}\label{eq:comp-weak-two}
  \Pstar\{\Lambda(p,Z)>t\}
  =\frac1n\E\#\{z:p(z)=0,\ \Lambda(p,z)>t\}
  \le \min\{1,Ct^{-2}\}.
\end{equation}
For every $0<q<2$ there is a constant $C_q$, independent of $n$, with
\[
  \Estar\Lambda^q\le C_q,
  \qquad 0\le\Estar\log\Lambda\le C.
\]
The moment range is sharp for bounds uniform in the degree:
\begin{equation}\label{eq:comp-second-asymptotic}
  \Estar\Lambda^2=2\log n+O(1)
  \qquad(n\to\infty).
\end{equation}
For $n=1$, the overlap is identically one.
\end{theorem}

The tail controls the expected fraction of poorly conditioned
eigenvalues.  It does not require independence of the roots and does not
assert a degree-independent bound for the worst root of a polynomial.
The second-moment assertion also shows that the exponent $2$ in
\eqref{eq:comp-weak-two} cannot be replaced by any larger exponent with
a degree-independent constant.

\subsection{An exact conditional moment}

Write $k_j(r)=\sum_{i=0}^{j}r^i$ and
$h_n(r)=(r k_n'(r)/k_n(r))'$, as in the box model.  At a fixed location
$z$, coefficient multiplication by $x-z$ is the injective map
\[
  D_z:\C^n\longrightarrow\C^{n+1},
  \qquad \mathbf q\longmapsto\operatorname{coeff}((x-z)q(x)).
\]
Its Gram matrix $T_z=D_z^*D_z$ is tridiagonal, with diagonal $1+r$,
superdiagonal $-z$, and subdiagonal $-\overline z$, where $r=|z|^2$.
In particular it is positive definite, including at $z=0$.

\begin{proposition}[Conditional overlap identity]\label{prop:comp-exact}
Under the selected-root law, the conditional second moment is
\begin{equation}\label{eq:comp-exact}
  \Phi_n(r):=\Estar[\Lambda^2\mid Z=z]
  =\frac{k_{n-1}(r)}{k_n(r)h_n(r)}\tr(T_z^{-1})
  =\frac{k_{n-1}(r)A_n(r)}{k_n(r)^2h_n(r)},
\end{equation}
where
\begin{equation}\label{eq:comp-trace-poly}
  A_n(r)=\sum_{j=0}^{n-1}(j+1)(n-j)r^j.
\end{equation}
The identity defines a continuous version of the conditional moment.
It satisfies
\[
  \Phi_n(r)=\Phi_n(r^{-1})\quad(r>0),
  \qquad \Phi_n(0)=n,
  \qquad \Phi_n(1)=\frac{2n}{n+1}.
\]
For $n=1$, $\Phi_1(r)=1$ for all $r\ge0$.
\end{proposition}

\begin{proof}
Ordinary Gaussian conditioning on $p(z)=0$ restricts the standard
Gaussian coefficient vector to the hyperplane $\operatorname{im}D_z$.
Writing $p(x)=(x-z)q(x)$ therefore gives
\[
  \mathbf q\mid\{p(z)=0\}\sim\CN(0,T_z^{-1}).
\]
Indeed its density in these coordinates is proportional to
$\exp(-\mathbf q^*T_z\mathbf q)$; the Jacobian of the linear
parametrization is constant.  Gaussian regression also gives
\[
  \E[|p'(z)|^2\mid p(z)=0]=k_n(r)h_n(r).
\]
At a selected root, the conditional law is tilted by $|p'(z)|^2$.
Since $p'(z)=q(z)$ and
$\Lambda^2=k_{n-1}(r)\|\mathbf q\|_2^2/|p'(z)|^2$, the tilt cancels
the denominator.  Thus
\[
  \Estar[\Lambda^2\mid Z=z]
  =\frac{k_{n-1}(r)\E[\|\mathbf q\|_2^2\mid p(z)=0]}
  {\E[|p'(z)|^2\mid p(z)=0]},
\]
which proves the first equality in \eqref{eq:comp-exact}.

The determinant of the $j\times j$ leading block of $T_z$ is $k_j(r)$:
both satisfy the recurrence
$d_j=(1+r)d_{j-1}-r d_{j-2}$, with $d_0=1$ and $d_1=1+r$.
Deleting its $i$th row and column leaves two such blocks, so
\begin{equation}\label{eq:comp-trace-minors}
  \tr(T_z^{-1})
  =\frac1{k_n(r)}\sum_{i=1}^{n}k_{i-1}(r)k_{n-i}(r)
  =\frac{A_n(r)}{k_n(r)}.
\end{equation}
The coefficient of $r^j$ in the numerator is $(j+1)(n-j)$,
which proves the second equality.

Reversal provides a useful check and removes the distinction between
inner and outer roots.  For $z\ne0$, let $p^{\mathrm{rev}}(x)=x^np(1/x)$.
Its coefficients have the same law as those of $p$, and direct reversal
of $\ell_z$ gives
\[
  \|\ell^{\mathrm{rev}}_{1/z}\|_2
   =|z|^{n-1}\|\ell_z\|_2,
  \qquad
  \|v(1/z)\|_2=|z|^{-(n-1)}\|v(z)\|_2.
\]
Hence $\Lambda(p^{\mathrm{rev}},1/z)=\Lambda(p,z)$.  Reversal carries a
uniformly selected root to a uniformly selected root, proving the
symmetry of $\Phi_n$.  Finally,
\[
  h_n(1)=\frac{n(n+2)}{12},
  \qquad \tr(T_1^{-1})=\frac{n(n+2)}6.
\]
Substitution proves the values at $r=1$; the assertions at $r=0$ and
$n=1$ follow directly from \eqref{eq:comp-exact}.
\end{proof}

\subsection{Uniform estimates and a tail bound}

For $0\le r\le1$, the cofactor formula
\eqref{eq:comp-trace-minors} immediately gives
\begin{equation}\label{eq:comp-trace-bound}
  \tr(T_z^{-1})
  \le n\min\{n,(1-r)^{-1}\}
  \le\frac{2n}{1-r+1/n},
\end{equation}
with the usual interpretation at $r=1$.  Combining this with
the metric estimate~\eqref{eq:box-metric-bound}, $k_{n-1}\le k_n$, and reversal proves
\begin{equation}\label{eq:comp-depth-bound}
  \Phi_n(r)\le C\bigl(1+n(1-\min\{r,r^{-1}\})\bigr)
  \le C(1+|n\log r|),\qquad r>0.
\end{equation}
The conditional second moment alone gives all the uniform moments of
order below two.  A little more of the same Gaussian regression gives
the sharper tail in Theorem~\ref{thm:comp-main}.

\begin{lemma}[Conditional moments]\label{lem:comp-third}
For each $0<q<4$ there is a constant $C_q$, independent of $n$ and $z$, such that
\[
  \Estar[\Lambda^q\mid Z=z]\le C_q\Phi_n(|z|^2)^{q/2}.
\]
\end{lemma}

\begin{proof}
Under ordinary conditioning on $p(z)=0$, put $\Sigma=T_z^{-1}$ and
$\sigma^2=v(z)^T\Sigma\overline{v(z)}=k_n(r)h_n(r)$.  Gaussian
regression decomposes the coefficient vector as
\[
  \mathbf q=b\xi+\eta,
  \qquad b=\Sigma\overline{v(z)}/\sigma,
  \qquad \xi=q(z)/\sigma\sim\CN(0,1),
\]
where $\eta$ is independent of $\xi$ and has covariance
$R=\Sigma-bb^*$.  At a selected root, the tilt is $|\xi|^2$.
It preserves the law of $\eta$ and its independence from $\xi$, while
$|\xi|^2$ now has density $u e^{-u}$ on $(0,\infty)$.  In particular
\[
  \Estar|\xi|^{-q}
    =\int_0^\infty u^{1-q/2}e^{-u}\,du
    =\Gamma(2-q/2)<\infty,\qquad 0<q<4.
\]
Writing $\beta=\|v(z)\|_2/\sigma$, we have
\[
  \Lambda=\beta\frac{\|b\xi+\eta\|_2}{|\xi|}
  \le\beta\|b\|_2+\beta\frac{\|\eta\|_2}{|\xi|},
  \qquad \Phi_n(r)=\beta^2\tr\Sigma.
\]
The standard Gaussian fourth-moment identity, valid also for a
degenerate covariance, is
\[
  \E\|\eta\|_2^4=(\tr R)^2+\tr(R^2)\le2(\tr R)^2.
\]
It follows by diagonalizing $R$ and using $\E|g|^2=1$ and
$\E|g|^4=2$ for a standard complex Gaussian $g$.
Hence $\E\|\eta\|_2^q\le2^{q/4}(\tr R)^{q/2}$ for $0<q<4$.  Apply
$(a+b)^q\le2^{\max\{q-1,0\}}(a^q+b^q)$, independence, and
$\|b\|_2^2\le\tr\Sigma$, $\tr R\le\tr\Sigma$ to obtain the claim.
\end{proof}

\begin{proposition}[Fixed-degree integrability]\label{prop:comp-fixed-degree}
For each fixed $n\ge2$ and $q>0$,
\[
  \Estar\Lambda^q<\infty\quad\Longleftrightarrow\quad q<4.
\]
For $n=1$ every moment is finite, since $\Lambda=1$.
\end{proposition}

\begin{proof}
The stronger first bound in \eqref{eq:comp-depth-bound} gives
$\Phi_n(r)\le Cn$ at every radius.  Lemma~\ref{lem:comp-third} therefore
proves finiteness for $q<4$.  For the converse, the residual covariance
$R$ in its proof has rank $n-1$: regression removes one complex
direction from the positive definite covariance $\Sigma$.  If $n\ge2$,
there is thus an $a>0$ with $\Prob\{\|\eta\|_2>a\}>0$.  On this event,
for sufficiently small $|\xi|$, one has $\|b\xi+\eta\|_2>a/2$.
Independence and the density $u e^{-u}$ of $|\xi|^2$ then show that
$\Estar[\Lambda^q\mid Z=z]=\infty$ for $q\ge4$, because
$\int_0^\varepsilon u^{1-q/2}e^{-u}\,du=\infty$.  Averaging over $Z$
proves the assertion.
\end{proof}

Thus the threshold four for each fixed degree comes from the small
derivative direction in the root-conditioned Gaussian.  The lower
threshold two for bounds uniform in degree also reflects the
distribution of root locations.

\begin{proof}[Proof of the tail and moment assertions in
Theorem~\ref{thm:comp-main}]
Let $D=|n\log|Z|^2|$ and $W=1+D$.  The box root law~\eqref{eq:box-depth} gives
$\Pstar\{D>u\}\le\min\{1,2/u\}$, so
$\Pstar\{W>u\}\le\min\{1,4/u\}$.  By conditional Markov,
Lemma~\ref{lem:comp-third} with $q=3$, and \eqref{eq:comp-depth-bound},
\[
  \Pstar\{\Lambda>t\}
  \le\Estar\min\{1,CW^{3/2}/t^3\}.
\]
For every $A>0$, integration of the tail gives
\[
  \Estar\min\{1,(W/A)^{3/2}\}
  =\frac{3}{2A^{3/2}}\int_0^A
       u^{1/2}\Pstar\{W>u\}\,du
  \le\min\{1,C/A\}.
\]
Taking $A$ to be a constant multiple of $t^2$ proves
\eqref{eq:comp-weak-two}.  Tail integration then proves
$\sup_n\Estar\Lambda^q<\infty$ for $0<q<2$ and
$\sup_n\Estar\log\Lambda<\infty$.  Nonnegativity of the logarithm
follows from $\Lambda\ge1$.
\end{proof}

\subsection{The critical second moment}

\begin{proposition}\label{prop:comp-second}
For every $n\ge1$,
\begin{equation}\label{eq:comp-second-integral}
  \Estar\Lambda^2
  =\frac2n\int_0^1
       \frac{k_{n-1}(r)A_n(r)}{k_n(r)^2}\,dr.
\end{equation}
In particular $\Estar\Lambda^2=2\log n+O(1)$ as $n\to\infty$.
\end{proposition}

\begin{proof}
The density of $|Z|^2$ is $h_n(r)/n$.  It cancels the $h_n$ in
\eqref{eq:comp-exact}, and reversal makes the contributions of
$(0,1)$ and $(1,\infty)$ equal.  This proves
\eqref{eq:comp-second-integral}, including finiteness for each $n$.

For $0\le r<1$, summing \eqref{eq:comp-trace-poly} gives
\[
  A_n(r)=
  \frac{n-(n+2)r+(n+2)r^{n+1}-nr^{n+2}}{(1-r)^3}.
\]
If $G_n(r)=k_{n-1}(r)A_n(r)/(n k_n(r)^2)$, rearrangement yields
\begin{align*}
  G_n(r)
  &=\frac1{1-r}
    \frac{(1-r^n)(1+r^{n+1})}{(1-r^{n+1})^2}\\
  &\hspace{1em}
    -\frac{2r}{n(1-r)^2}
      \left(\frac{1-r^n}{1-r^{n+1}}\right)^2.
\end{align*}
Suppose $n\ge2$.  On $0\le r\le1-1/n$ we have
$r^{n+1}\le e^{-1}$, and therefore
\[
  \left|G_n(r)-\frac1{1-r}\right|
  \le\frac{Cr^n}{1-r}+\frac{2}{n(1-r)^2}.
\]
The right side has a uniformly bounded integral on this interval:
for its first term use $r^n\le e^{-n(1-r)}$ and
$u=n(1-r)\ge1$; the second term integrates directly.  Hence
\[
  \int_0^{1-1/n}G_n(r)\,dr=\log n+O(1).
\]
On the remaining interval, \eqref{eq:comp-trace-bound} and
$k_{n-1}\le k_n$ give $0\le G_n(r)\le n$, so its integral is at most
one.  Multiplication by two proves the asserted asymptotic.
\end{proof}

This completes Theorem~\ref{thm:comp-main}.  By monotonicity of moments
on a probability space, $\sup_n\Estar\Lambda^q=\infty$ for every
$q\ge2$, allowing the value $+\infty$ for a fixed degree.  Thus the
sharpness assertion concerns uniformity in $n$; it does not say that
all moments of order at least two diverge at each fixed degree.

\begin{example}[Degrees one and two]\label{ex:comp-small}
When $n=1$, $C_p$ is a scalar and $\Lambda=1$.  When $n=2$,
\[
  T_z=\begin{pmatrix}1+r&-z\\-\overline z&1+r\end{pmatrix},
  \qquad
  \Phi_2(r)=\frac{2(1+r)^2}{1+4r+r^2}.
\]
In particular the conditional second moment is $4/3$ on the unit
circle, and direct integration of \eqref{eq:comp-second-integral}
gives
\[
  \Estar\Lambda^2=\frac23+\frac{4\pi}{9\sqrt3}.
\]
The deterministic overlap can nevertheless become arbitrarily large:
if $p(x)=(x-z)(x-w)$ with $z\ne w$, then
$\Lambda(p,z)=\sqrt{(1+|z|^2)(1+|w|^2)}/|z-w|$.
The theorem describes how averaging over the polynomial and a selected
root controls these near-collision configurations.
\end{example}

\subsection{Comparison with coefficient conditioning}

For the same scalar output $z$, coefficient and matrix perturbations
give, respectively,
\[
  \kappa_{\mathrm{coeff}}(p,z)
     =\frac{\sqrt{k_n(|z|^2)}}{|p'(z)|},
  \qquad
  \kappa_{\mathrm{eig}}(C_p,z)=\Lambda(p,z).
\]
All $n+1$ polynomial coefficients are perturbable in the first
condition number, including the leading coefficient.  Multiplication
by the corresponding data norm and division by the common output norm
$|z|$ define their relative versions; $z\ne0$ almost surely.

\begin{corollary}[A relative conditioning comparison]\label{cor:comp-relative}
With Euclidean coefficient norm and spectral matrix norm, define
\[
  \kappa_{\mathrm{coeff,rel}}(p,z)
    =\frac{\sqrt{k_n(|z|^2)}\|\mathbf a\|_2}{|p'(z)|\,|z|},
  \qquad
  \kappa_{\mathrm{eig,rel}}(C_p,z)
    =\frac{\Lambda(p,z)\|C_p\|_2}{|z|}.
\]
Then
\begin{align*}
  \Estar\log\frac{\kappa_{\mathrm{eig}}}{\kappa_{\mathrm{coeff}}}
    &=\log n+O(1),\\
  \Estar\log\frac{\kappa_{\mathrm{eig,rel}}}
                          {\kappa_{\mathrm{coeff,rel}}}
    &=\log n+O(1).
\end{align*}
The error terms are bounded uniformly in the degree.
\end{corollary}

\begin{proof}
The normalized derivative
$p'(Z)/\sqrt{k_n(|Z|^2)h_n(|Z|^2)}$ has the determinant-weighted
complex Gaussian law from the local theorem: its squared modulus has
density $u e^{-u}$.  Consequently the following finite-degree identity
holds, where $\gamma$ is Euler's constant:
\begin{equation}\label{eq:comp-absolute-identity}
  \Estar\log\frac{\kappa_{\mathrm{eig}}}{\kappa_{\mathrm{coeff}}}
  =\Estar\log\Lambda
    +\frac12\Estar\log h_n(|Z|^2)+\frac{1-\gamma}{2}.
\end{equation}
The bounded logarithmic overlap and the box-model estimate~\eqref{eq:box-metric-log}
$\Estar\log h_n(|Z|^2)=2\log n+O(1)$ prove the first assertion.

For the relative comparison the output factors cancel exactly.  It
remains to control
$\E\log\|C_p\|_2-\E\log\|\mathbf a\|_2$; a statistic depending only
on the polynomial has the same expectation under $\Pstar$, because
there are always $n$ roots.  With
$b=(a_0,\ldots,a_{n-1})^T/a_n$, the last column gives the lower bound
below, while $C_pC_p^*=\diag(0,1,\ldots,1)+bb^*$ gives the upper bound:
\[
  \|b\|_2\le\|C_p\|_2\le\sqrt{1+\|b\|_2^2}.
\]
Since sums of $m$ independent squared standard complex Gaussians have
the $\operatorname{Gamma}(m,1)$ law, their logarithmic means give
\begin{equation}\label{eq:comp-data-scale}
  \frac\gamma2-\frac1{2n}
  \le \E\log\|C_p\|_2-\E\log\|\mathbf a\|_2
  \le\frac\gamma2.
\end{equation}
For example, this follows from
$\E\log\operatorname{Gamma}(m,1)=H_{m-1}-\gamma$, with $H_0=0$.
Adding \eqref{eq:comp-data-scale} proves the relative assertion.
\end{proof}

These are signed logarithmic comparisons of two specified local
condition numbers.  They say that the geometric mean of their ratio is
bounded above and below by constant multiples of $n$.  The comparison
does not include errors in forming $C_p$ or in recovering information
from an eigenvector; an end-to-end stability theorem would also have to
analyze those maps.

\section{Interpretation and further questions}\label{sec:discussion}

The local and companion results describe two different sources of sensitivity.
The normalized Jacobian law is determined by a first-order Gaussian calculation
and the determinant weight at a root. Its inverse moments have threshold four.
The companion overlap additionally contains the coefficient norm of a Lagrange
polynomial. For a fixed degree $n\ge2$ its moment threshold is again four, but degree
uniformity has threshold two. The radial estimates and the exact second moment
show why concentration of roots near the unit circle does not permit discarding
the complementary radii from conditioning averages.

This distinction also explains the relation to the limiting spectral measures
in \cite{barbarino2023monic,barbarino2022limit}. At a simple eigenvalue $\lambda$
of a matrix $M$, the spectral projector satisfies
\[
 \Pi_\lambda=\lim_{z\to\lambda}(z-\lambda)(zI-M)^{-1}
 =\frac{rw^*}{w^*r},\qquad \|\Pi_\lambda\|_2
 =\frac{\|r\|_2\|w\|_2}{|w^*r|}.
\]
Consequently an overlap estimate requires control of a resolvent residue.
Barbarino--Noferini's estimates for shifted monic companions
\cite[Lemma~A.11]{barbarino2023monic} concern fixed spectral parameters away
from the unit circle; they cannot simply be evaluated at a random eigenvalue,
where the shifted matrix is singular. Their decomposition into a deterministic
block shift and a Gaussian correction, followed by the Woodbury identity,
nevertheless suggests a useful route for studying such residues in block
companion models.

\subsection{A finite-degree illustration}

Table~\ref{tab:second-moments} evaluates the exact integral
\eqref{eq:comp-second-integral}. The conditional second moment on the unit
circle stays below two, while the unconditional second moment grows. These
values illustrate the proved formulas; they are not used to infer either
asymptotic or a uniform bound.

\begin{table}[htbp]
\centering
% Generated by research/paper-ac/second_moment_table.py
\begin{tabular}{rrrr}
\toprule
$n$ & $\Phi_n(1)$ & $\Estar\Lambda^2$ & $\Estar\Lambda^2-2\log n$\\
\midrule
1 & 1.000000 & 1.000000 & 1.000000 \\
2 & 1.333333 & 1.472800 & 0.086505 \\
8 & 1.777778 & 2.849863 & -1.309020 \\
32 & 1.939394 & 4.876708 & -2.054764 \\
128 & 1.984496 & 7.342059 & -2.362002 \\
\bottomrule
\end{tabular}

\caption{Deterministic quadrature of the exact second-moment formula.
Positive-coefficient Horner evaluation avoids cancellation at $r=1$.
Computations at 30 and 50 decimal digits agreed to more than 25 decimal places;
the table reports six. Degrees one and two also agree with their exact values
in Example~\ref{ex:comp-small}.}
\label{tab:second-moments}
\end{table}

The calculation is reproducible from
\texttt{research/paper-ac/second\_moment\_table.py}. A separate symbolic script
checks the cofactor polynomial, reversal, the conditional variance, and the
closed second-moment integrand in small degrees. These computations are
consistency checks of the identities, rather than replacements for the proofs.

\subsection{What a conditioning comparison establishes}

Corollary~\ref{cor:comp-relative} says that the geometric mean of the relative
companion-to-coefficient condition ratio is comparable to $n$. The ratio uses
the same scalar output but two different input spaces. It is compatible with
bounded average logarithmic eigenvector overlap because the Euclidean
coefficient-to-root condition has its own degree scale. Neither statement
alone bounds floating-point error in a complete algorithm.

An analysis of a particular implementation would have to control formation of
the companion matrix, the structure of the backward error from the eigensolver,
and any subsequent recovery. Equation~\eqref{eq:setup-composition} explains
the role of structure: a perturbation induced exactly by polynomial
coefficients has the original root derivative, whereas an arbitrary matrix
perturbation need not arise from a comparably small coefficient perturbation.
Other linearizations can change this comparison
\cite{toh1994pseudozeros,deteran2017eigenvalue}.

The selected-root tail also gives, by a union bound,
\[
 \Prob\{\max_{p(z)=0}\Lambda(p,z)>t\}\le\min\{1,Cn/t^2\}.
\]
The factor $n$ is part of this bound. Controlling a typical root, an expected
fraction of roots, and the worst root are different probabilistic objectives.

\subsection{The multivariate problem left by the local theory}

For systems, the representation question is now isolated in
$\Theta_{n,d}^{\B}$ from \eqref{eq:box-decomposition}. The local law and the
box metric estimates already hold in every dimension. What remains is to
control interpolation angles for an explicitly specified quotient basis or
algebraic construction. The univariate division map $D_z$ has a particularly
simple tridiagonal Gram matrix; a multivariate replacement must account for
relations between polynomial multiples and for basis selection.

Two extensions require additional ideas. First, bounds for a fixed bivariate
construction would begin to connect the local theory with multivariate
algebraic rootfinding. Any dependence on the dimension must be retained if
the objective is to compare with dimension-exponential worst cases. Second,
an additive perturbation model $P=P_0+\sigma Q$, where $P_0$ is fixed and $Q$
is a centered Gaussian system, produces noncentral conditional
derivatives and a noncentral determinant weight. The centered independence
law of Theorem~\ref{thm:local-universal} cannot be transferred to that smoothed
model without further analysis. The results proved here provide the centered
benchmark for these questions. Smoothed local conditioning and homotopy
complexity are already studied in the homogeneous Bombieri--Weyl setting
\cite{burgisser2011smale,armentano2016condition}; see also
Section~\ref{sec:local-classical}. The remaining task here is to control the
specified coefficient and representation conditions for the box ensemble
and its noncentered perturbations.

The noncentral Gaussian estimate of
\cite[Theorem~A.5]{barbarino2023monic} is a possible ingredient in this task,
not a substitute for the determinant-weighted law at a selected root.
Moreover, matrices produced by multivariate elimination generally reuse the
same input coefficients in several entries. Applying random matrix-polynomial
results to them requires treating these dependencies explicitly.


\bibliographystyle{plainurl}
\bibliography{gaussian-conditioning}
\end{document}
